%%%%%%%%%%%%%%%%%%%%%%%%%%%%%%%%%%%%%%%%%%%%%%%%%%%%%%%%%%%%%%%%%%%%%%%%%%%%%%%%%%
\begin{frame}[fragile]\frametitle{}
\begin{center}
{\Large ध्यान --- शिवकृपानंद स्वामी}
\end{center}

{\tiny (Ref: शिवकृपानंद स्वामी (समर्पण ध्यान) --- ध्यान दिवस मुलाखत, YouTube)}

\end{frame}

%%%%%%%%%%%%%%%%%%%%%%%%%%%%%%%%%%%%%%%%%%%%%%%%%%%%%%%%%%%
\begin{frame}[fragile]\frametitle{ध्यान म्हणजे काही न करणं}

	\begin{itemize}
	\item ध्यान म्हणजे: \textbf{काही न करणं} --- शरीराने नाही, विचाराने नाही
	\item `ध्यान करणे' हा शब्दच चुकीचा --- ध्यान \textbf{लागते}, केले जात नाही
	\item ध्यान म्हणजे आपली \textbf{मूळ स्वाभाविक स्थिती} --- प्रकृतीशी समरस होणे
	\item जंगलात एकटे फिरा, समुद्रकाठी एकटे फिरा --- विचार आपोआप येत नाहीत
	\item आजच्या पिढीचे चित्त कमकुवत --- TV, मोबाईल, इंटरनेटमुळे बाहेर गेलेले
	\item चित्त अंतर्मुखी करणारा एकमात्र: \textbf{गुरु}
	\end{itemize}

{\tiny (Ref: शिवकृपानंद स्वामी --- ध्यान)}

\end{frame}

%%%%%%%%%%%%%%%%%%%%%%%%%%%%%%%%%%%%%%%%%%%%%%%%%%%%%%%%%%%
\begin{frame}[fragile]\frametitle{गुरु तुम्हाला निवडतो}

	\begin{itemize}
	\item आपण गुरुला निवडू शकत नाही --- \textbf{गुरु तुम्हाला निवडतो}
	\item रामदास स्वामी: `तुमची आध्यात्मिक बैठक चांगली करा, गुरु घरी येऊन भेटेल'
	\item या गुरु चांगले का ते गुरु --- तुलना करत बसू नका; कुठेतरी बसा, पुढे जा
	\item गुरु आवश्यक आहे --- आत्मसाक्षात्कारासाठी गुरुशिवाय मार्ग नाही
	\item जिवंत गुरु महत्त्वाचा: `बिना बीज वृक्ष उगत नाही' --- समाधिस्थ गुरु प्रेरणा, अनुभूती नाही
	\item गुरुतत्त्व एक आहे --- रूपे वेगळी, परंतु सर्वांमध्ये गुरुतत्त्व एकच
	\end{itemize}

{\tiny (Ref: शिवकृपानंद स्वामी --- ध्यान)}

\end{frame}

%%%%%%%%%%%%%%%%%%%%%%%%%%%%%%%%%%%%%%%%%%%%%%%%%%%%%%%%%%%
\begin{frame}[fragile]\frametitle{आध्यात्मिक बीज --- गुरुशिवाय प्रगती नाही}

	\begin{itemize}
	\item शेतकरी रोज पाणी-खत टाकतो --- पण आंब्याचे बीजच नाही तर झाड कसे?
	\item जप, पूजापाठ, कर्मकांड --- बीजच नाही तर आध्यात्मिक प्रगती नाही
	\item \textbf{आध्यात्मिक बीज} = गुरूकडून मिळणारी अनुभूती
	\item ती अनुभूती मिळाल्यावर आत्मसाक्षात्कार --- आणि मग वाईट कर्म घडणारच नाही
	\item वाईट कर्म करण्यासाठी वाईट शक्ती लागते --- अनुभूतीनंतर ती शक्तीच उरत नाही
	\item ध्यानाच्या प्रवाहातूनच ती आत्मअनुभूती मिळू शकते
	\end{itemize}

{\tiny (Ref: शिवकृपानंद स्वामी --- ध्यान)}

\end{frame}

%%%%%%%%%%%%%%%%%%%%%%%%%%%%%%%%%%%%%%%%%%%%%%%%%%%%%%%%%%%
\begin{frame}[fragile]\frametitle{शिवकृपानंद यांचा आध्यात्मिक प्रवास}

	\begin{itemize}
	\item नागपूरचे --- आजीकडून बालपणी संस्कार: `देव अविनाशी आहे'
	\item पाच वर्षांचे असताना: `कोणतेही शरीर देव असू शकत नाही' --- हे कळले
	\item मुसलमान फकीर, हिंदू गुरु --- जात-धर्म न पाहता सर्वांना आशीर्वाद --- हे पाहिले
	\item एम.कॉम. + मार्केटिंग नोकरी --- पण वाटायचे `हे माझे काम नाही'
	\item नदीकाठी बसताना विचार: `हे पाणी समुद्राला मिळेल तसे मी परमात्म्यात केव्हा मिळेन?'
	\item \textbf{तीव्र तळमळ निर्माण झाल्यावरच गुरूचे आगमन होते}
	\end{itemize}

{\tiny (Ref: शिवकृपानंद स्वामी --- ध्यान)}

\end{frame}

%%%%%%%%%%%%%%%%%%%%%%%%%%%%%%%%%%%%%%%%%%%%%%%%%%%%%%%%%%%
\begin{frame}[fragile]\frametitle{नेपाळमधील गुरुभेट --- तीन दिवसांची समाधी}

	\begin{itemize}
	\item ध्यानात सतत दिसायचे: शिव मंदिर, पशुपतीनाथ, एक उंच बाबाजी
	\item कानपूरहून काठमांडूला गेले --- तिथे एक वृद्ध तीन दिवसापासून वाट पाहत होते
	\item नेपाळच्या शिबू गावाजवळ गुफेत राहणारे `शिवबा' --- शिवाचा साक्षात अवतार
	\item गुफेतून बाहेर येताना ओळखले: `हेच माझ्या ध्यानात येतात'
	\item शिवबा म्हणाले: `४० वर्षांपासून या वाट पाहतोय'
	\item गुफेत डोक्यावर हात ठेवला --- \textbf{तीन दिवस समाधी} लागली
	\item समाधीनंतर: `देवाचा शोध संपला; देव माझ्यामध्येच आहे'
	\end{itemize}

{\tiny (Ref: शिवकृपानंद स्वामी --- ध्यान)}

\end{frame}

%%%%%%%%%%%%%%%%%%%%%%%%%%%%%%%%%%%%%%%%%%%%%%%%%%%%%%%%%%%
\begin{frame}[fragile]\frametitle{आत्मानुभूती --- देव तुमच्यातच आहे}

	\begin{itemize}
	\item जुन्या मंदिरात दर्शनासाठी सहा तास रांग --- एक क्षण दर्शन --- पुन्हा जायला तयार
	\item का? कारण त्या एक क्षणात \textbf{आत्मानुभूती} झालेली असते
	\item देव मूर्तीत नाही --- \textbf{देव तुमच्यामध्येच आहे} --- ते एक क्षण जाणवते
	\item ती अनुभूती एक क्षणाची --- ध्यान साधनेने \textbf{कायमस्वरूपी} होते
	\item करोना काळात मंदिरे बंद --- साधकांना फरक नाही: `देव आमच्यातच आहे'
	\item स्वामींच्या आठ आश्रमांत एकही देवाचे मंदिर नाही
	\end{itemize}

{\tiny (Ref: शिवकृपानंद स्वामी --- ध्यान)}

\end{frame}

%%%%%%%%%%%%%%%%%%%%%%%%%%%%%%%%%%%%%%%%%%%%%%%%%%%%%%%%%%%
\begin{frame}[fragile]\frametitle{ओरा / अभामंडळ --- चित्त शुद्धीचे प्रतिबिंब}

	\begin{itemize}
	\item ओरा म्हणजे तुमचे चित्त किती शुद्ध, किती पवित्र आहे ते दिसते
	\item कर्करोग होण्याच्या सहा महिने आधी विचारांत कर्करोग असतो --- योगी ओळखतो
	\item वैद्यकीय सत्यापन: अहमदाबाद मेडिकल कॉलेज --- १५०० डॉक्टरांच्या ओरा तपासणीत सिद्ध
	\item तीन दिवस ध्यान शिबिरानंतर ओऱ्यात स्पष्ट बदल दिसला
	\item चित्त नाशवान गोष्टींवर टाकल्यास चित्त कमकुवत होते
	\item आकाश पहा --- नाशवान नाही; तिकडे चित्त टाकल्यास बळकट होते
	\end{itemize}

{\tiny (Ref: शिवकृपानंद स्वामी --- ध्यान)}

\end{frame}

%%%%%%%%%%%%%%%%%%%%%%%%%%%%%%%%%%%%%%%%%%%%%%%%%%%%%%%%%%%
\begin{frame}[fragile]\frametitle{ध्यानातील तीन अडथळे}

	\begin{itemize}
	\item \textbf{पहिला हमला --- शरीर:} सवय नसल्याने विरोध करते
		\begin{itemize}
		\item खाज सुटते, जांभई येते, गॅसेस होतात, खोकला येतो --- सर्व शरीराचे प्रकार
		\end{itemize}
	\item \textbf{दुसरा हमला --- विचार:} ५२ वर्षापूर्वी घडलेल्या क्षुल्लक गोष्टी आठवतात
		\begin{itemize}
		\item चित्त शुद्ध होताना साचलेले विचार बाहेर पडतात --- नाली ढवळल्यासारखे
		\end{itemize}
	\item \textbf{तिसरा हमला --- बुद्धी:} `तुझ्याने ध्यान होणार नाही, महत्त्वाचे काम सोडून बसलास'
	\end{itemize}

{\tiny (Ref: शिवकृपानंद स्वामी --- ध्यान)}

\end{frame}

%%%%%%%%%%%%%%%%%%%%%%%%%%%%%%%%%%%%%%%%%%%%%%%%%%%%%%%%%%%
\begin{frame}[fragile]\frametitle{ट्रॅफिक पोलीस उपाय}

	\begin{itemize}
	\item ट्रॅफिक पोलीस: ट्रक, बस, स्कूटर येतात --- पाहतो, काही करत नाही
	\item तसे विचारांकडे \textbf{फक्त पाहायचे} --- त्यांत सहभागी व्हायचे नाही
	\item विचारात ऊर्जा दिली की त्यामागे आणखी विचार येतात
	\item पहिला विचार आला: `ऑफिसचा विचार, ठीक आहे' --- पुढे जाऊ द्यायचे
	\item पहिला प्याला रोकता आला तर व्यसन होत नाही --- पहिला विचार रोका
	\item बुद्धीचे काम: उठवणे --- तिचा उद्देश एकच, त्याला ओळखा
	\end{itemize}

{\tiny (Ref: शिवकृपानंद स्वामी --- ध्यान)}

\end{frame}

%%%%%%%%%%%%%%%%%%%%%%%%%%%%%%%%%%%%%%%%%%%%%%%%%%%%%%%%%%%
\begin{frame}[fragile]\frametitle{व्यसनमुक्ती आणि ध्यान}

	\begin{itemize}
	\item नशामुक्ती केंद्रात: डॉक्टरांनी सांगितले `लिव्हर सडेल' --- मान हलवले, सवय गेली नाही
	\item स्वामींनी म्हटले: `दारू सोडवायला आलो नाही, परमात्म्याशी जोडायला आलो'
	\item ४५ दिवसांचे अनुष्ठान केले --- \textbf{९०\% लोकांची दारू सुटली}
	\item कारण: ध्यानाने आत्मभाव वाढतो --- सर्व व्यसने आपोआप सुटतात
	\item नाभचक्र विकसित झाल्यावर शरीर चांगले नाही ते आपोआप बाहेर फेकते
	\item `दोन फुटांची आत्मभावाची रेषा खेच --- बाकी लहान होऊन जाते'
	\end{itemize}

{\tiny (Ref: शिवकृपानंद स्वामी --- ध्यान)}

\end{frame}

%%%%%%%%%%%%%%%%%%%%%%%%%%%%%%%%%%%%%%%%%%%%%%%%%%%%%%%%%%%
\begin{frame}[fragile]\frametitle{४५ दिवसांचे अनुष्ठान}

	\begin{itemize}
	\item \textbf{अर्धा तास गुरूला दान द्या} --- दान केलेल्यावर अधिकार नाही
	\item ध्यान लागणे-न लागणे तुमचे क्षेत्र नाही --- फक्त अर्धा तास बसणे तुमचे काम
	\item रोज सकाळी + संध्याकाळी अर्धा तास --- ४५ दिवस सलग
	\item आधी सराव: एका \textbf{निर्जन स्थानी} अर्धा तास बसा
		\begin{itemize}
		\item मोबाईल, पेपर, पुस्तक नाही; गाणे, शिट्टी नाही; शेजारी कोणी नाही
		\item यामुळे विचार हळूहळू कमी होतात
		\end{itemize}
	\item शरीर मोबाईलसारखे --- आठ तास एकदा नको, अर्धा तास रोज पुरे
	\item एकही साधक असा नाही की ज्याने ४५ दिवस केले आणि आयुष्य बदलले नाही
	\end{itemize}

{\tiny (Ref: शिवकृपानंद स्वामी --- ध्यान)}

\end{frame}

%%%%%%%%%%%%%%%%%%%%%%%%%%%%%%%%%%%%%%%%%%%%%%%%%%%%%%%%%%%
\begin{frame}[fragile]\frametitle{शुद्ध आत्मा भावना + आत्मसाक्षात्कार प्रार्थना}

	\begin{itemize}
	\item उजव्या हाताचा तळवा डोक्यावर ठेवा (टाळू भागावर), तीनदा दाबा, घड्याळाच्या दिशेने फिरवा
	\item डोळे बंद करा, हात खाली घ्या; चित्त टाळू भागावर ठेवा
	\item मनात म्हणा:
		\begin{itemize}
		\item \textit{मैं एक पवित्र आत्मा हूं} $\cdot$ \textit{याने शरीर} $\cdot$ \textit{एक शुद्ध आत्मा}
		\item \textit{शुद्ध आत्मा केवल परमात्मा} $\cdot$ \textit{परमात्मा के सिवा मेरा कोई अस्तित्व नहीं}
		\end{itemize}
	\item प्रार्थना: \textit{हे परमात्मा, कृपया मुझे आत्मसाक्षात्कार प्रदान करे}
	\item ध्यानाचा उद्देश: \textbf{शरीरभाव संपावा, आत्मभाव जागृत व्हावा}
	\end{itemize}

{\tiny (Ref: शिवकृपानंद स्वामी --- ध्यान)}

\end{frame}
